\documentclass[11pt]{article}

\usepackage[a4paper, margin=1in]{geometry}
\usepackage[T1]{fontenc}
\usepackage[utf8]{inputenc}
\usepackage{amsmath, amssymb}
\usepackage{booktabs}
\usepackage{graphicx}
\usepackage{hyperref}
\usepackage{xcolor}
\hypersetup{colorlinks=true, linkcolor=blue, urlcolor=blue, citecolor=blue}
\usepackage{authblk}
\usepackage{microtype}
\usepackage{enumitem}
\setlist[itemize]{leftmargin=1.4em}
\setlist[enumerate]{leftmargin=1.6em}

\newcommand{\grad}{\nabla}
\newcommand{\half}{\tfrac{1}{2}}

\title{\textbf{A Proton--Electron Cosmology}\\[0.25em]
\large RealQM and Newtonian Gravitation from a Single Electric-Potential Seed}

\author[1]{Claes Johnson}
\affil[1]{KTH Royal Institute of Technology, Stockholm, Sweden \\ \texttt{claesjohnson@gmail.com}}
\date{\today}

\begin{document}
\maketitle

\begin{abstract}
We sketch \textbf{RealUniv}, an exploratory cosmology built from the RealQM / RealNucleus program, in which all matter is Coulombic
(charge clouds in ordinary Euclidean space, no strong force, no elementary neutron) and all gravitation is
Newtonian. The novelty here is a \textbf{single primitive seed}: a small, small-scale fluctuation of the electric
\emph{potential} $\varphi_E$ alone --- charge being its Laplacian, $\rho_c=-\varepsilon_0\grad^2\varphi_E$. From it
the universe develops in two tiers. \textbf{At the small scale}, the Laplacian inflates the potential seed into
charge of both signs and of \emph{different spatial extent} --- small dense positive (protons) and large diffuse negative
(electrons); since a contracted cloud carries more energy, that extent \emph{is} the inertial mass ($E/c^2$), and
Coulomb geometry then builds nuclei (RealNucleus) and atoms (RealQM). \textbf{At the large scale}, the
\emph{energy} of that charged matter is a mass/inertia density $\rho$, and it is this $\rho$ --- seeded entirely by
the charges --- that sources gravitation through Poisson, $\rho=\grad^2\varphi_g$. Read about the mean, $\rho$
carries both signs, so the gravitational source is two-fold; the perturbations separate into great chunks of
positive mass \emph{interfoliated} with negative mass --- the cosmic web of filaments and clusters interleaved with
under-dense voids, the negative sector acting as a repulsive dark energy. There is thus \emph{no primitive negative
mass and no separate gravitational seed}: only the electric potential is seeded, mass is the energy of the charged world, and
gravity is emergent. The construction is strictly non-relativistic --- Coulomb below, Newton above --- so the speed
of light $c$ enters only the electromagnetic sector: there is a local speed limit for charges and light, but none
for the large-scale gravitational motion of neutral bodies (two galaxies may approach at relative speed $>c$). We
state the picture, note that it \emph{reproduces} the primordial helium fraction as a consistency (with the
observed $n$--$p$ gap $E_\ast$ and the standard freeze-out supplied, not derived), and set out its load-bearing
open problems --- chiefly \emph{what sets the nuclear scale}, the same open datum as $E_\ast$.

\medskip
\noindent\textbf{Status.} An exploratory proposal, not established cosmology. It departs from the Standard Model and
$\Lambda$CDM and must eventually be reconciled with them; the open problems of Section~\ref{sec:open} are
load-bearing. Read it as ``what follows \emph{if}.''

\medskip
\noindent\textbf{Keywords:} cosmology; two-tier cosmology; proton--electron model of the nucleus; emergent
gravitation; dark energy; primordial nucleosynthesis; RealQM.
\end{abstract}

\section{Introduction}
Cosmology is unavoidably speculative at its foundations. The standard account --- hot Big Bang, cosmic
inflation, $\Lambda$CDM --- fits the data only by invoking ingredients that are inferred rather than seen: an
initial singularity, an inflationary epoch, dark matter, and dark energy, none of them observed in a laboratory,
each introduced to close a gap. In a subject resting on such foundations it is
legitimate, and useful, to ask what a \emph{simple, rational, non-standard} alternative would look like --- one
built not from new fields and new epochs but from the most ordinary ingredients physics already has:
\textbf{electron--proton Coulomb interaction and Newtonian gravitation}. This note develops one such model.

The RealQM program formulates quantum mechanics as multiphase 3D continuum mechanics --- unit-charge densities on
non-overlapping domains separated by free boundaries, minimising the Coulomb energy~\cite{realqm}. Its nuclear
extension, RealNucleus, revives the pre-1932 proton--electron picture (neutron $=$ proton $+$ electron) and
reproduces nuclear binding from pure Coulomb interactions, with no strong-force postulate~\cite{realnucleus}. This
note asks what cosmology results if the \emph{same} ingredients are taken as the starting point of the universe,
together with ordinary Newtonian gravitation.

The answer, which we call \textbf{RealUniv}, is a \textbf{two-tier} cosmology from \textbf{one seed}. The seed is a
fluctuation of the electric \emph{potential} (charge being its Laplacian). It
shapes the small scale directly (charge $\to$ nuclei $\to$ atoms, all of RealQM); and the \emph{energy} of that
charged matter, being mass, sources an \emph{emergent} Newtonian gravitation that shapes the large scale (the
cosmic web, with a repulsive negative-mass sector as dark energy). We keep strictly to what Coulomb and Newton can
say, mark the extreme regimes (strong gravity, the highest energies) where more would be needed, and keep the open
problems explicit.

\section{The machinery: two-valued Poisson, opposite force-signs}\label{sec:machinery}
Both forces share \emph{one} structure and differ in a single sign. Each is the gradient of a potential obeying
Poisson's equation with a two-valued ($\pm$) source,
\begin{equation}
\text{electric:}\quad \rho_c=-\varepsilon_0\,\grad^2\varphi_E,
\qquad\qquad
\text{gravitational:}\quad \rho_m=\frac{1}{4\pi G}\,\grad^2\varphi_g,
\end{equation}
the force being $-(\text{charge})\,\grad\varphi$ in each case. In this cosmology the electric potential is the
primitive seed (Section~\ref{sec:seed}) and the gravitational source is the emergent energy of the charged world
(Section~\ref{sec:large}) --- but the \emph{machinery} is common. The one structural difference is the relative
sign, and with it the force law:
\begin{center}
\begin{tabular}{lll}
\toprule
 & like sources & unlike sources \\
\midrule
electric (charge) & \textbf{repel} & attract \\
gravitational (mass) & \textbf{attract} & repel \\
\bottomrule
\end{tabular}
\end{center}
Concretely there are two independent binary signs --- the source $\pm\grad^2\varphi$ (two-valued) and the force
$\mathbf F=\pm\grad\varphi$ (repulsive or attractive for like sources) --- and their \emph{four} combinations are
exactly the four ``charges'' the two forces couple to: the repulsive-for-like sign is electromagnetism, its $\pm$
source the $\pm$charge of the small scale; the attractive-for-like sign is gravitation, its $\pm$ source the $\pm$
mass/energy of the large scale. One operator and two signs thus yield $\{\pm\text{charge},\,\pm\text{mass}\}$.

This single sign difference generates the two tiers. Because like charges \emph{repel} and unlike attract, matter
neutralises into $\pm$ pairs, charge cancels, and electromagnetism \emph{screens} --- it stays local and
small-scale. Because like masses \emph{attract}, mass \emph{clumps} and accumulates, growing without bound into
large-scale structure; and because unlike masses repel, positive- and negative-mass regions \emph{separate} into
the interfoliated web and voids of Section~\ref{sec:large}. The scale separation is thus a \emph{consequence} of
the force-sign, not a separate assumption: small-scale electromagnetism and large-scale gravitation are one
two-valued Poisson structure, read with two signs. (This is the static Poisson level; the electromagnetic sector
carries $c$ and radiation, and $\pm$-mass gravitational dynamics has its own subtleties --- both belonging to the
extreme regimes outside the Coulomb--Newton scope kept here.)

A second asymmetry --- in the \emph{direction} the Poisson relation is used --- fixes the \emph{scale} of each
force. For electromagnetism the \emph{potential} is primary: the seed is a small, \emph{small-scale} (short-
wavelength) oscillation of $\varphi_E$ --- not a smooth ripple --- and charge is obtained by \emph{differentiating},
$\rho_c=-\varepsilon_0\grad^2\varphi_E$. The Laplacian is a \emph{local} operator that \emph{magnifies} short
wavelengths ($\rho_k\propto k^2\varphi_k$), so precisely because the seed is small-scale (large $k$) it is inflated,
despite its tiny amplitude, into sharp, small-scale $\pm$ charges. For gravitation the \emph{source} is primary --- the mass/energy of
the charged world --- and the potential is obtained by \emph{integrating}, $\varphi_g=4\pi G\,\grad^{-2}\rho_m$.
The inverse Laplacian is a \emph{global} operator (convolution with the $1/r$ Green's function) that \emph{smooths},
suppressing short wavelengths ($\varphi_k\propto\rho_k/k^2$) and building a broad, large-scale potential. So the
same Poisson relation, run forward makes small-scale electromagnetism and run backward makes large-scale
gravitation: \textbf{the Laplacian acts locally and magnifies; the inverse Laplacian acts globally and smooths} ---
and that is why charge is small and gravity large.

\section{The arena and the single seed}\label{sec:seed}
The arena is a plain \textbf{3D Euclidean coordinate space} --- absolute space, simply given, with no physical
medium. The only operator needed is the Laplacian $\grad^2$, and the
only primitive is a small-scale fluctuation of the \textbf{electric potential} $\varphi_E$. Charge is not
postulated separately; it is the curvature of that potential,
\begin{equation}
\rho_{\text{charge}} = -\varepsilon_0\,\grad^2\varphi_E ,
\end{equation}
with two immediate consequences. \textbf{Both signs:} $\rho_{\text{charge}}$ is positive where $\varphi_E$ is
concave, negative where convex. \textbf{Net zero:} $\int\grad^2\varphi_E\,dV$ is a surface term that vanishes, so
the universe carries \textbf{zero net charge} automatically --- equal numbers of protons and electrons, with no
baryogenesis. And because the Laplacian multiplies each Fourier mode by $k^2$ ($\rho_k\propto k^2\varphi_k$), a
small, \emph{small-scale} (short-wavelength) oscillation of $\varphi_E$ --- not a smooth ripple --- becomes a large
density contrast: the seed is tiny in amplitude but short in wavelength, and it is the short wavelength (large $k$)
that the Laplacian amplifies into strong charge.

Crucially, \textbf{only the electric potential is seeded} (charge being its Laplacian). Mass and gravitation are not primitive; they arise, in
Section~\ref{sec:large}, as the \emph{energy} of the charged structures the seed produces. This is the essential
departure from a cosmology of two independent potentials: here there is one seed, one potential, and gravity is a
consequence.

\section{The small scale: RealQM and RealNucleus}\label{sec:small}
The seed, inflated by the Laplacian, produces charge of both signs and, decisively, of \emph{different spatial extent}: small dense
positive charge and large diffuse negative charge. These are the \textbf{proton} and the \textbf{electron}, and
the whole charged world is built from them by Coulomb geometry.

\paragraph{Size is inertial mass.} In RealQM a particle is a charge cloud whose size is fixed by the balance of a
localisation (kinetic) cost against Coulomb attraction. A \emph{contracted} cloud costs more localisation and
Coulomb energy and so carries \emph{more energy}; and inertia is energy content, $\text{mass}=E/c^2$~\cite{nomass}.
Hence \textbf{small $=$ heavy, large $=$ light}: the proton (small, dense) is heavy, the electron (large, diffuse)
is light. ``Proton small, electron big'' and ``proton heavy, electron light'' are one fact, not two --- a size
difference read through $E/c^2$. (The numerical ratio $m_p/m_e=1836$ is not derived; it is the open input of
Section~\ref{sec:open}, as it is in every theory.)

\paragraph{No overlap, no self-interaction: self-repulsion is kinetic energy.} RealQM charge clouds do not overlap
--- each electron occupies its own domain with a free boundary --- and carry \emph{no self-interaction}: there is
no Coulomb self-repulsion term. The role such self-repulsion would play, resisting the collapse of a charge to a
point, is carried instead by the \emph{localisation energy} $\half\int|\grad\psi|^2$: compressing a cloud sharpens
its gradients and costs kinetic energy, so it resists compression and settles at a finite size (localisation cost
against Coulomb attraction). A free electron, unconfined, therefore spreads; a bound one takes the size where
localisation balances attraction. In RealQM, \textbf{self-repulsion appears as kinetic energy}.

\paragraph{Nuclei and atoms, one mechanism at two scales.} With a neutron taken as $p+e$, a neutral atom $(A,Z)$
holds $A$ protons and $A$ electrons: $A-Z$ \emph{compact} inside the nucleus and $Z$ \emph{expanded} in orbitals.
The electron occupies two size-phases, its extent adapting to the Coulomb well: compact in a deep well (inside a
nucleus, ${\sim}$fm, ${\sim}$MeV) and expanded in a shallow one (atomic, ${\sim}$\AA, ${\sim}$eV). Chemistry and
nuclear physics are then one Coulomb free-boundary mechanism at two scales: a molecule is nuclei glued by expanded
(valence) electrons; a nucleus is protons glued by compact electrons. RealNucleus reproduces nuclear binding this
way --- the He\,-4/deuteron binding ratio comes out $13.1$ against an experimental $12.7$, from Coulomb geometry
alone~\cite{realnucleus}, and light-nucleus fusion (hydrogen burning) appears as a rearrangement of protons and
electrons, not the action of a separate force.

\paragraph{The nuclear/atomic scale.} What sets the two sizes can be read from charge and the proton's finite size
$R_p$ alone, with \emph{no mass and no $c$}. An electron attracted to one proton in vacuum must taper to zero at
infinity; the localisation cost then forces a large cloud of size the atomic $a_0=\hbar^2/me^2\approx0.5$\,\AA. An
electron caged among protons at a multiphase free boundary need \emph{not} taper to zero --- its density hands off
to the proton density --- so a flat caged density has zero localisation cost, and it sits at the cage (proton) size
$R_p\approx$ a few fm, bound by Coulomb (${\sim}$MeV at fm). The size ratio is therefore
\begin{equation}
\frac{a_0}{R_p}\;\approx\;\frac{1}{\alpha^2}\;\approx\;2\times10^{4},
\end{equation}
so the fine-structure constant appears here as a pure \emph{geometric} ratio of the two scales,
$\alpha^2=R_p/a_0$, the ratio of nuclear to atomic size --- not a coupling constant. Empirically the H atom
($\approx a_0=0.53$\,\AA) and the deuteron nucleus ($\approx 2.1$\,fm) give a ratio $\approx1/25{,}000$, matching
$\alpha^2\approx1/18{,}800$ to order of magnitude. \emph{Why} the proton is ${\sim}$fm --- what fixes $R_p$ ---
remains the central open problem.

\section{Inertia is energy; the local speed limit}\label{sec:inertia}
Throughout, inertia is the energy content: $\mathbf p=(E/c^2)\mathbf v$, the measured relation of the Penning
trap~\cite{nomass}, and Newton's second law $\mathbf F=\dot{\mathbf p}$ governs the response. In the small, the
force is electromagnetic and couples to \emph{charge}; in the large, it is gravity and couples to \emph{mass}
$=E/c^2$. One mechanics, one inertia, two forces reading two properties of the same matter.

The constant $c$ enters \emph{only} the electromagnetic sector --- through the field momentum of a moving charge,
Thomson's electromagnetic mass. Consequently there is a \textbf{local speed limit}: a charge, or a pulse of light,
cannot overtake light \emph{through} the local matter and fields, because feeding it energy only increases its
inertia $E/c^2$ without bound as $v\to c$. But \textbf{Newtonian gravitation is $c$-free} --- instantaneous, with
no $c$ in $\rho=\grad^2\varphi_g$ --- so it imposes \emph{no} speed limit on the large-scale motion of neutral
bodies. Two galaxies may approach, or recede, at relative speed \emph{greater than $c$}; this is a global relative
velocity of gravitating matter, not a local signal, and it violates no local causality (no local frame sees
anything pass light through it). The speed limit is thus a property of the RealQM/Maxwell sector, not of gravity ---
just as in observational cosmology distant galaxies recede faster than light without any local law being broken.
(How Newtonian mechanics extends to relatively moving frames --- where $c$ and its kinematics would enter --- is a
separate matter, taken up in \emph{Many-Minds Relativity}~\cite{mmr}, and lies outside the non-extreme scope kept
here.)

\section{The large scale: emergent Newtonian gravitation}\label{sec:large}
The energy of the charged matter --- protons, electrons, their binding and radiation --- is a mass/inertia density
$\rho$, \textbf{positive throughout} ($\rho\ge 0$), seeded entirely by the charges. It is this positive $\rho$ that
sources gravitation, and the emergent potential $\varphi_g$ comes in \textbf{two forms}, its Laplacian taking
either sign against the one positive source:
\begin{equation}
\pm\,\frac{1}{4\pi G}\,\grad^2\varphi_g \;=\; \rho \;\ge\; 0 .
\end{equation}
Thus $\varphi_g$ is not a primitive but \emph{emergent}, the large-scale field of the aggregated energy of the
charged world, and it comes in two forms \textbf{according to the sign of its Laplacian}: the $+$ form is a
\emph{well} ($\grad^2\varphi_g>0$, effective positive mass, attractive), the $-$ form an \emph{anti-well}
($\grad^2\varphi_g<0$, effective negative mass, repulsive). The $\pm$ mass is therefore \emph{formal} --- the sign
carried by the curvature of $\varphi_g$, not a negative energy; the source $\rho$ is positive. By the gravitational force-sign, the wells
attract and draw together while the anti-wells repel, so the positive-mass regions are \textbf{isolated} into great
chunks --- filaments and clusters --- \emph{interfoliated} with the anti-well regions between them: the
\textbf{cosmic web}, positive-mass nodes and filaments interleaved with voids, the anti-well (formally
negative-mass) sector acting as \textbf{dark energy}.

\emph{Both} the positive and the negative mass are emergent --- the two forms of the gravitational potential of the
one charge-seeded, positive energy $\rho$ --- neither a primitive, neither a second seed; they repel and
\textbf{separate}, the positive isolating into the web and the negative emptying into the voids. Only charge is put
in; the whole mass--gravity structure, of both signs, is the emergent large-scale field of that energy. This is the
sense in which the largest structures and the smallest share one origin: the small scale is the seed and its charge
directly, the large scale the emergent gravity of their energy.

\paragraph{A 2D test: cosmic-web morphology.} A mass-conserving 2D simulation of a $\pm$mass field
$\rho=\grad^2\varphi$ from a random potential, evolved as compressible self-gravitating flow with same-sign
attraction and opposite-sign repulsion, develops the morphology of observed large-scale structure
(Figure~\ref{fig:web}): intermixed positive- and negative-mass filamentary webs, mass concentrated at nodes,
separated by large voids --- the voids here \emph{actively driven} by negative-mass repulsion rather than merely
evacuated by gravity. A qualitative match (a 2D toy); quantitative tests (void-size distribution, correlation
function) remain to be done.
\begin{figure}[h]\centering
\includegraphics[width=0.5\textwidth]{cosmic_web_pm2.png}
\caption{2D self-gravitating $\pm$mass from $\rho=\grad^2\varphi$: positive- (warm) and negative- (cool) mass webs
with concentrated nodes and large voids --- cosmic-web morphology, voids driven by negative-mass repulsion.
Mass-conserving; periodic box.}\label{fig:web}
\end{figure}

\section{Dark energy and inflation from one seed}\label{sec:de-infl}
Two central features of modern cosmology appear here as consequences of the \emph{same} machinery, from the one
charge seed --- an economy worth stating, and honestly bounding.

\emph{Dark energy.} The under-dense, formally negative-mass sector of Section~\ref{sec:large} is repulsive: the
anti-wells push matter away, driving the voids to grow and separate. This is a built-in \textbf{dark-energy
mechanism}, not put in by hand --- it is simply the $-$ side of the density perturbation, the anti-well form of the
emergent $\varphi_g$. What is \emph{not} yet derived is its magnitude, or that its effective equation of state is
$w\approx-1$; the mechanism is qualitative.

\emph{Inflation.} The Laplacian amplifies short wavelengths, $\rho_k\propto k^2\varphi_k$, so a tiny small-scale
oscillation of the potential is magnified into a large density contrast --- an \textbf{``inflation'' of the density
field}. This is amplification of the \emph{field}, not the exponential expansion of \emph{space} of the standard
inflationary paradigm (which addresses the horizon and flatness problems); the two must not be conflated.

Because $\rho=\grad^2\varphi$ is $k^2$-weighted, the amplification is short-wavelength dominated: the seed comes
out as \emph{small-scale} charge structure. In this model that is not a spectral defect to be observed in the sky
but the very mechanism by which the initial fluctuation becomes a \textbf{proton--electron universe} --- nuclei,
atoms, bound matter at small scales (Section~\ref{sec:small}). One might object that a $k^2$-weighted (``blue'')
density spectrum conflicts with the near scale-invariant, slightly red primordial spectrum ($n_s\approx0.96$)
inferred from the CMB. But that inference is internal to the hot Big Bang: it reads the CMB as the red-shifted
\emph{afterglow} of a hot dense phase and reconstructs a primordial density spectrum through recombination and a
transfer function. This model has no Big Bang and no afterglow phase; it makes no claim about the CMB and requires
none. The $n_s$ comparison therefore presupposes a framework the model does not adopt, and the short-wavelength
weighting is read here as a feature --- the origin of small-scale matter --- rather than a conflict. The
interpretation of the CMB itself we regard as a separate, open observational question outside the scope of a
small-scale-structure cosmology.

So from a single seed the model contains natural \emph{candidates} for both dark energy (negative mass) and a
fluctuation amplification (``inflation of the density field'') --- an economy that a separate-mechanism cosmology
does not have --- while neither is yet quantitatively validated.

\section{The cosmic sequence and the helium fraction}\label{sec:sequence}
The two sectors exchange energy in \emph{both} directions, and the cosmic sequence is that exchange in time.
\emph{Upward}, the energy of the charged world is mass and sources gravity (Section~\ref{sec:large}).
\emph{Downward}, gravitational attraction feeds energy back into electromagnetism: the collapse of a positive-mass
region converts gravitational potential energy into heat, and that hot dense state drives the electromagnetic
processes. Gravity supplies the energy; electromagnetism spends it as structure and light.

Concretely: ordinary gravitational contraction of a positive-mass seed supplies one hot dense epoch, and the
released binding energy leaves as light. As the contracting proton--electron matter heats, \textbf{deuterons and nuclei} form
(${\sim}$MeV, RealNucleus), the Coulomb binding energy radiating away (the origin of light) and appearing as the
mass defect. On cooling, nuclei capture expanded electrons into \textbf{atoms} (${\sim}$eV, RealQM), again
radiating; at that recombination the matter turns neutral, electromagnetism screens out, and radiation streams
free. Only \emph{then} does the weak, slow gravitational aggregation build the cosmic web of
Section~\ref{sec:large}. So the universe's timeline stages the two tiers in order: \emph{charge builds the small
while matter is charged; gravity builds the large once it is neutral}.

A bound state survives once the bath can no longer ionise it, so the two electron phases switch on at two
temperatures set by their two binding energies, forcing the order plasma $\to$ nuclei $\to$ atoms
(Table~\ref{tab:epochs}). The temperature ratio between the epochs (${\sim}10^6$) is just the energy-scale ratio.

\begin{table}[h]\centering
\begin{tabular}{lll}
\toprule
epoch & process & temperature \\
\midrule
plasma & free $\pm$ charges & $>10^{10}$ K \\
nuclei & compact $e$ captured $\to$ nuclei & ${\sim}10^{10}$ K ($kT{\sim}$MeV) \\
atoms & expanded $e$ captured $\to$ atoms; radiation & ${\sim}3000$ K ($kT{\sim}$eV) \\
\bottomrule
\end{tabular}
\caption{Two condensation epochs, separated by ${\sim}10^6$ in temperature.}\label{tab:epochs}
\end{table}

\paragraph{The helium fraction: a consistency, not a prediction.} The proton--electron picture reproduces the
primordial helium; it is worth being exact about what is derived and what is borrowed. The electron has two
size-phases --- \emph{compact} (inside a nucleus, i.e.\ inside a neutron, since $\text{neutron}=p+\text{compact
}e$) and \emph{expanded} (atomic) --- interconverted by the two weak processes, compactify ($p+e\to n$) and expand
($n\to p+e$, $\beta$-decay). Three quantities enter, and all three are \emph{inputs}:
\begin{itemize}
\item $E_\ast$ --- the \textbf{energy gap} between the two phases: the energy to compactify an electron, released
when it expands. It plays the role the neutron--proton mass difference plays in standard cosmology and is
\emph{identified with it}, $E_\ast\approx1.29$~MeV. It is an \emph{input}, not a RealNucleus output: its absolute
size is the open nuclear-scale problem (Section~\ref{sec:open}) --- RealNucleus fixes dimensionless \emph{ratios},
not the MeV scale.
\item $T_\ast$ --- the \textbf{transition temperature}, $T_\ast\sim E_\ast/k\sim10^{10}$~K: above it the two phases
interconvert freely, below it compactification becomes rare.
\item $T_f$ --- the \textbf{freeze-out temperature} ($\approx0.8$~MeV): where the weak (size-changing) processes
fall slower than the cosmic expansion and the neutron/proton ratio \emph{locks}. As in standard big-bang
nucleosynthesis it is set by the weak-rate-vs-expansion balance, not predicted --- a second input.
\end{itemize}
In equilibrium the neutron/proton ratio is the Boltzmann factor $e^{-E_\ast/kT}$; at freeze-out it locks at
$(n/p)_f=e^{-E_\ast/kT_f}$; a fraction $\beta$-decay back (survival factor $d$), and essentially all the rest end in
He-4:
\begin{equation}
\Big(\frac{n}{p}\Big)_{\text{nuc}}=\frac{(n/p)_f\,d}{1+(n/p)_f(1-d)},
\qquad
Y=\frac{2\,(n/p)_{\text{nuc}}}{1+(n/p)_{\text{nuc}}} .
\end{equation}
With $E_\ast\approx1.29$~MeV, $T_f\approx0.8$~MeV and $d\approx0.74$: $(n/p)_f=0.199\to(n/p)_{\text{nuc}}=0.140\to
Y=0.245$, the observed primordial helium. This is a \textbf{consistency}, not a new prediction: with the two inputs
borrowed --- $E_\ast$ the observed $n$--$p$ gap, $T_f$ the standard freeze-out --- the size-transition picture
\emph{reproduces} the standard result. Nothing here is derived; and, notably, RealNucleus's one genuine computed
output, the He-4/deuteron binding ratio ($13.1$ vs $12.7$), does \emph{not} enter $Y$. The content is that the
proton--electron reinterpretation is \emph{consistent} with the standard helium calculation, with $E_\ast$ the same
open datum as the nuclear scale.

\section{Complexity from the size asymmetry}\label{sec:complexity}
The world is built not on a charge imbalance (charge is exactly balanced, $\pm e$) but on a \emph{size} imbalance:
the proton has a fixed small size (a rigid anchor) while the electron's size is variable, adapting to its well ---
${\sim}$\AA{} in an atom, ${\sim}$fm in a nuclear cage. The fixed proton supplies the anchors; the variable
electron \emph{spans scales} --- spreading to bind atoms and molecules, compacting to bind nuclei --- producing the
nested multi-scale hierarchy that \emph{is} complexity. With a single fixed size only one scale exists and nothing
complex is built. The nuclear/atomic separation is thus a \textbf{precondition for complexity}, so any universe
that supports structure, and hence observers, must exhibit it: an anthropic reason \emph{why} the scales are
separated, though not a derivation of \emph{how far}. Notably, since a neutron is structurally $p+e$, it is heavier
than a proton \emph{by construction}, so proton and hydrogen stability is automatic --- the usual fine-tuning
bound ``$\alpha$ must not be too large or the proton decays'' simply does not arise.

\section{Relation to $\Lambda$CDM and observations}\label{sec:lcdm}
It is worth stating plainly what this model shares with the standard picture, where it differs, and what it does
not yet address.

\emph{Shared.} Ordinary matter is Coulombic and gravitation is attractive between masses; nuclei and atoms form and
later assemble into a filamentary large-scale structure of clusters and voids; a helium mass fraction near
$Y\approx0.25$ is accommodated.

\emph{Different.} There is no initial singularity and no hot Big Bang, hence no cosmic microwave background as a
red-shifted afterglow (Section~\ref{sec:de-infl}); no inflationary expansion of space (the amplification here is of
the density \emph{field}, not of space); no dark-matter particle (structure is sourced by the ordinary
energy/mass density of Coulombic matter); and no cosmological constant (a repulsive dark-energy sector arises as
the negative-curvature part of the emergent gravitational potential). The neutron is composite ($p+e$), not
elementary, and nuclear binding is Coulombic geometry rather than a strong force.

\emph{Not yet addressed.} The model is deliberately restricted to Coulomb$+$Newton and does not, as it stands,
account for the redshift--distance relation and the expansion history, the CMB acoustic-peak spectrum under any
alternative interpretation, precise large-scale-structure statistics (the two-point correlation and its baryon
acoustic feature), or light-element abundances beyond helium. These are the observational fronts on which any
serious successor must eventually be tested; here we claim only internal consistency and economy of ingredients,
not empirical parity with $\Lambda$CDM. What the model offers in exchange is fewer free primitives --- one seed, no
dark-matter species, no separate inflaton, no cosmological constant --- at the cost of predictive machinery still
to be built (Section~\ref{sec:open}).

\section{Open problems}\label{sec:open}
\begin{enumerate}
\item \textbf{The nuclear scale.} What sets the ${\sim}10^4$--$10^5$ nuclear/atomic separation --- equivalently,
the proton size $R_p$? The preferred reading takes $R_p$ (${\sim}$fm) as the input nuclear length, with the caged,
zero-localisation electron sitting at it (no relativity needed); the ratio $a_0/R_p\approx1/\alpha^2$ then follows.
Why $R_p$ has that value is open. This is the same datum as the mass ratio $m_p/m_e$ and the proton/electron size
ratio: one open number.
\item \textbf{The fluctuation spectrum and a law of motion.} Here $\rho=\grad^2\varphi$ is a definition, not yet a
dynamics: $\varphi$ needs its own law of motion, and the seed spectrum is an input. The $k^2$-weighting makes the
amplification short-wavelength dominated, which this model reads as the origin of small-scale matter rather than as
a primordial-spectrum prediction (Section~\ref{sec:de-infl}); comparison with the CMB-inferred $n_s\approx0.96$
presupposes the Big Bang afterglow interpretation, which the model does not adopt. The status of the CMB in a
non-Big-Bang setting is itself an open question, but not one this small-scale-structure model needs to settle.
\item \textbf{The mass ratio $m_p/m_e=1836$.} Built in qualitatively as a size difference, but its value is not
derived (nor is it in the Standard Model).
\item \textbf{The neutrino.} $\beta$-decay's continuous spectrum --- the 1930 objection that sank the
proton--electron neutron~\cite{pauli}. The ``electron expands'' step needs a home for that energy and spin.
\item \textbf{Precision and scope.} BBN abundances beyond helium, the fluctuation-to-structure statistics, and the
extreme regimes (strong-field gravity, the highest energies) lie beyond the Coulomb$+$Newton scope kept here.
\end{enumerate}

\section{Summary}\label{sec:summary}
\noindent On a given 3D Euclidean space, if the sole primitive is a small-scale fluctuation
of the electric \emph{potential} (charge being its Laplacian) --- yielding small heavy protons and large light
electrons, hence nuclei and atoms by Coulomb geometry
(RealQM/RealNucleus) --- then the \emph{energy} of that charged matter, single quantity, sources an emergent
Newtonian gravitation whose $\pm$ perturbations separate into the cosmic web (positive mass) and dark energy
(negative mass), with the speed of light limiting only the electromagnetic sector; a unifying two-tier picture from
one seed, whose survival hinges on a single hard question: what sets the nuclear scale.

\begin{thebibliography}{9}
\bibitem{realqm} C.\ Johnson, \emph{RealQM: Quantum Mechanics as Multiphase 3D Continuum Mechanics}, and
interactive gallery, \url{https://claes542.github.io/RealMolecule/}.
\bibitem{realnucleus} C.\ Johnson, \emph{RealNucleus: nuclear binding from a pure-Coulomb proton--electron model},
manuscript, same repository.
\bibitem{nomass} C.\ Johnson, \emph{Energy Without Mass: RealQM, Inertia, and the Penning Trap}, same repository.
\bibitem{mmr} C.\ Johnson, \emph{Many-Minds Relativity}: absolute Euclidean space, $c$ as a wave speed, and a
critique of Einstein's special relativity.
\bibitem{pauli} W.\ Pauli, open letter to the Tübingen group (1930), proposing the neutrino to save energy
conservation in $\beta$-decay.
\bibitem{kolbturner} E.\ W.\ Kolb and M.\ S.\ Turner, \emph{The Early Universe}, Addison-Wesley (1990).
\end{thebibliography}

\end{document}
